\Article{365}
Uns genügt es hier, die Möglichkeit dieser Lösungen gezeigt zu haben. Viele Kunstgriffe, die diese Operationen erleichtern können, übergehen wir der Kürze wegen. Einige sehr wichtige Folgerungen können wir aber nicht auslassen.

Durch das Vorige ist bewiesen, dass alle Hilfsgleichungen zur Lösung der Gleichung $x^\nu=1$, wenn man sie in Kongruenzen umwandelt, mögliche Wurzeln haben, wenn die Periode
\SourceLeaf{244}
\[
x+x^p+x^{pp}+\dotstwo+x^{p^{m-1}}
\]
noch nicht zerlegt ist. Bleiben wir bei dem Fall, dass $\nu$ eine Primzahl ist; dann ist $m$ ein Teiler von $\nu-1$. Hier haben also die Hilfskongruenzen wirkliche Wurzeln, wenn die Anzahl der durch sie gefundenen Perioden ein Teiler von $\frac{\nu-1}m$ ist. Ist daher $\frac{\nu-1}m$ gerade, ist also $m$ ein Teiler von $\frac{\nu-1}2$, oder gilt gleichbedeutend $p^{\frac{\nu-1}2}\equiv1\;(\mathrm{mod.}\,\nu)$, oder ist gleichbedeutend $p$ ein quadratischer Rest der Primzahl $\nu$, so hat die quadratische Gleichung, durch die die Wurzeln in zwei Perioden geteilt werden, wirkliche Wurzeln modulo $p$. In Kapitel \textsc{vi} haben wir aber gezeigt, dass diese Gleichung, wenn man $\nu=4n\pm1$ setzt, stets $xx+x\mp n=0$ ist. Somit erhalten wir den bemerkenswerten

\textsc{Satz.} \emph{Ist die Primzahl $p$ ein quadratischer Rest der Primzahl $4n\pm1$, so hat die Kongruenz}
\[
xx+x\mp n\equiv0(\mathrm{mod.}\,p)
\]
\emph{wirkliche Wurzeln und folglich auch die Kongruenz}
\[
4xx+4x\mp4n\equiv0\quad\text{oder}\quad(2x+1)^2\mp\nu\equiv0
\]
\emph{das heißt, $\pm\nu$ ist ein quadratischer Rest von $p$.}

\Article{366}
Dies ist also der dritte vollständige Beweis des Fundamentalsatzes aus Kapitel \textsc{iv}. Er verdient umso mehr Aufmerksamkeit, als die Prinzipien, aus denen er gewonnen ist, von denen, die wir für die früheren Beweise verwendet haben, völlig verschieden sind.\ .Aus derselben Quelle wollen wir nun auf dem entgegengesetzten Weg einen vierten ableiten. Es sei nämlich $\nu$ eine Primzahl der Form $4n\pm1$, $p$ eine beliebige andere Primzahl, und $\pm\nu$ sei ein quadratischer Rest der Primzahl $p$; wir werden beweisen, dass $p$ ein quadratischer Rest von $\nu$ ist.

\Needspace{34\baselineskip}Es sei $p^m$ die kleinste Potenz von $p$, die $\equiv1(\mathrm{mod.}\,\nu)$ ist. Die elementaren Teiler der Funktion $\frac{x^\nu-1}{x-1}$ modulo $p$ haben den Grad $m$, sodass ihre Gesamtzahl $=\frac{\nu-1}m$ ist. Da nun $\pm\nu Rp$ gilt, ist die Kongruenz
\[
xx+x\mp n\equiv0(\mathrm{mod.}\,p)
\]
lösbar; ihre Wurzeln seien $A,A'$. Man verteile die Funktionen $x,xx,\dotstwo x^{\nu-1}$ auf zwei Klassen, die mit $\xi,\xi'$ bezeichnet werden. Dann ist
\SourceLeaf{245}
\[
\begin{aligned}
\xi+\xi'&\equiv A+A'+(1+x+xx+\dotstwo+x^{\nu-1})\\
\xi\xi'&\equiv AA'+\lambda(1+x+xx+\dotstwo+x^{\nu-1})
\end{aligned}
\]
also ist
\[
(z-\xi)(z-\xi')-(z-A)(z-A')
\]
durch jeden elementaren Teiler der Funktion $\frac{x^\nu-1}{x-1}$ teilbar. Daher teilt jeder dieser elementaren Teiler entweder $\xi-A$ und $\xi'-A'$ oder $\xi-A'$ und $\xi'-A$. Hieraus folgt (da $A$ nicht $\equiv A'$ ist): Setzt man für $x$ nun $x^p$, so ändern sich $\xi$ und $\xi'$ nicht. Ginge nämlich $\xi$ in $\xi'$ über und umgekehrt, so wären $\xi-A$ und $\xi-A'$ durch dieselbe Primfunktion teilbar. Q.\ E.\ A. Schließlich folgt daraus, dass $\frac{\nu-1}2$ durch $m$ teilbar ist oder gleichbedeutend $p^{\frac{\nu-1}2}-1$ durch $\nu$. Also ist $p$ ein quadratischer Rest von $\nu$. Q.\ E.\ D.

Alle Fälle des Fundamentalsatzes lassen sich leicht aus den beiden Sätzen ableiten.

\Article{367}
Obwohl wir uns hier auf den Fall beschränkt haben, dass $\nu$ eine Primzahl ist, lassen sich auch für zusammengesetztes $\nu$ ohne große Mühe analoge Sätze bestimmen; der Kürze wegen können wir dies jetzt nicht ausführlicher darstellen.

Es ist offensichtlich, dass sich ähnliche Beobachtungen auch über eine größere Anzahl von Perioden anstellen lassen. Ist also $\frac{\nu-1}m$ durch $3$ teilbar, ist also $p$ ein kubischer Rest der Primzahl $\nu$, so ist die Gleichung, durch die die Wurzeln der Gleichung $x^\nu=1$ auf drei Perioden verteilt werden und deren Bestimmung a priori wir in Kapitel \textsc{vi} gelehrt haben, modulo $p$ lösbar und umgekehrt. So lässt sich zum Beispiel die Kongruenz $x^3+xx-2x-1\equiv0$ modulo jeder Primzahl der Form $7n\pm1$ lösen, nicht aber, wenn sie eine andere Form hat.

Es würde uns nicht schwerfallen, dieses Kapitel mit vielen anderen Beobachtungen zu bereichern, wenn die Grenzen, innerhalb deren wir bleiben müssen, es nicht verhinderten. Denjenigen, die weitergehen möchten, können diese Prinzipien wenigstens den Weg weisen.
